\section{Cyclotomic constraints and supported local conditions}
\label{sec:irreg-setup}

Finite torsors on the proper Fermat curve allow the classical first-case
arguments to use the local representatives of quotient sections.
Kummer's unit argument and its cyclotomic setting are presented in
\citet[Theorem~1.1 and Proposition~1.5, pp.~1--6]{Washington1982}.
Herbrand's class-group implication is
\citet[\S4, pp.~423--431]{Herbrand1932}; we use the direct character
convention recalled by
\citet[Theorem~1.1 and the following paragraph, p.~151]{Ribet1976}.
The constructions below establish the passage from those arithmetic
inputs to sections without assuming rational Fermat coordinates.

Throughout, \(p\geq5\) is prime. Let
\[
 F_p:\ X^p+Y^p=Z^p,\qquad
 D_p=F_p\cap\{XYZ=0\},\qquad U_p=F_p\setminus D_p.
\]
Here \(F_p\) is the minus-sign model of \(X_p\), identified with
the all-plus model by changing the sign of \(Z\); the notation
\(U_p\) is used through that identification. Finite fields are
written \(\F_p\). Its three rational coordinate cusps are
\[
 Q_0=(0:1:1),\qquad Q_1=(1:0:1),\qquad Q_\infty=(1:-1:0).
\]
Write \(J_p=\Jac(F_p)\) and \(j_p(R)=[R-Q_0]\).

Let \(\Delta_p=\pi_1(F_{p,\overline{\Q}})\). Define
\[
 \Pi_p^{\ab,p}
 =
 \pi_1(F_p/\Q)/
 \overline{[\Delta_p,\Delta_p]\Delta_p^p}.
\]
This quotient retains the full arithmetic quotient \(G_\Q\).
We consider a geometric quotient
\[
 \pi_1(U_p/\Q)\twoheadrightarrow\Pi\twoheadrightarrow G_\Q
\]
equipped with a compatible map \(\Pi\to\Pi_p^{\ab,p}\), and a continuous
section \(s:G_\Q\to\Pi\). At each finite place \(v\), suppose that its
restriction is, up to geometric conjugacy, induced by a full local
point section or by a full section of the decomposition group of a
\(\Q_v\)-rational point of \(D_p\). Such a representative has a proper
centre \(R_v\in F_p(\Q_v)\). The arguments below require only that the
image in \(\Pi_p^{\ab,p}\) be represented by the full proper point
section of \(R_v\).

All assertions concern every permitted choice of the \(R_v\). In
particular, uniqueness of a local point from a finite abelian
cohomology class is not assumed. A centre at \(p\) is called
\emph{first case} if its primitive integral homogeneous coordinates
are all units in \(\Z_p\).

Set
\[
 K=\Q(\zeta_p),\qquad L=\Q_p(\zeta_p),\qquad
 \Delta=\Gal(K/\Q),\qquad
 \sigma_i(\zeta_p)=\zeta_p^i,\quad
 \omega(\sigma_i)=i\in\F_p^\times.
\]
For an \(\F_p[\Delta]\)-module, a superscript \((\omega^m)\)
always denotes the direct eigenspace on which \(\sigma_i\) acts
as multiplication by \(i^m\).

\subsection{Proper Kummer torsors and virtual units}
\label{sec:irreg-torsors}

Define the virtual-unit space
\[
 V=\left\{[\alpha]\in K^\times/K^{\times p}:
             v(\alpha)\equiv0\pmod p
             \text{ for every finite place }v\right\}.
\]
Let \(E_K=\mathcal O_K^\times\), and let a bar denote complex conjugation.

\begin{lemma}\label{lem:irreg-virtual}
There is an exact sequence of \(\F_p[\Delta]\)-modules
\begin{equation}
 0\longrightarrow E_K/E_K^p
   \longrightarrow V
   \longrightarrow\Cl(K)[p]\longrightarrow0.
 \label{eq:irreg-virtual}
\end{equation}
Consequently, if \(r^-=\dim_{\F_p}\Cl(K)[p]^-\), then
\begin{equation}
 \dim_{\F_p}V^-=1+r^-.
 \label{eq:irreg-minus-dim}
\end{equation}
For odd \(m\) with \(3\leq m\leq p-2\),
\begin{equation}
 V^{(\omega^m)}\simeq\Cl(K)[p]^{(\omega^m)}.
 \label{eq:irreg-odd-components}
\end{equation}
\end{lemma}

\begin{proof}
If \((\alpha)=\mathfrak a^p\), send \([\alpha]\) to
\([\mathfrak a]\). Changing \(\alpha\) by a \(p\)-th power changes
\(\mathfrak a\) by a principal ideal. Every class killed by \(p\)
is obtained in this way, and the kernel consists of units times
\(p\)-th powers. This proves \eqref{eq:irreg-virtual}, including its
equivariance for the natural actions on elements and ideals.

For \(\epsilon\in E_K\), the unit
\(\epsilon/\overline\epsilon\) has absolute value one under every
complex embedding and is therefore a root of unity. The roots of
unity in \(K\) are \(\mu_{2p}\). Since \(2\) is invertible modulo \(p\),
the odd part of \(E_K/E_K^p\) is the line generated by \(\zeta_p\);
this line has character \(\omega\). It is nonzero, since a \(p\)-th
root of \(\zeta_p\) in \(K\) would be a root of unity of order divisible
by \(p^2\). Finally, \(\#\Delta=p-1\) is prime to \(p\), so taking
eigenspaces is exact. This gives both conclusions.
\end{proof}

Use signed coordinate functions \(A,B,C\), obtained by permuting
\(X,Y,-Z\), so that
\[
 A^p+B^p+C^p=0.
\]
For \(i\in\Z/p\Z\), let \(P_i\) be the point where \(C=0\) and
\(A+\zeta_p^iB=0\), and put
\[
 g_i=\frac{A+\zeta_p^iB}{A+B}.
\]

\begin{lemma}\label{lem:irreg-proper-torsors}
The section supplies classes
\[
 \xi_i\in H^1(K,\mu_p)=K^\times/K^{\times p}
\]
satisfying \(\xi_i=\sigma_i\xi_1\) for \(i\neq0\).
If its chosen \(p\)-local centre is first case, then every \(\xi_i\)
belongs to \(V\). If that centre has signed coordinates \(a,b,c\) and
\(e=b/(a+b)\), then
\begin{equation}
 \loc_{\mathfrak p}(\xi_i)
  =[1-e+e\zeta_p^i]\in L^\times/L^{\times p}.
 \label{eq:irreg-local-values}
\end{equation}
Here \(\mathfrak p\) is the unique prime of \(K\) above \(p\).
\end{lemma}

\begin{proof}
The tangent line \(A+\zeta_p^iB=0\) meets the Fermat curve in the
divisor \(pP_i\). Thus
\begin{equation}
 \divisor(g_i)=p(P_i-P_0).
 \label{eq:irreg-divisor}
\end{equation}
The equation \(W^p=g_i\) consequently extends to an unramified
\(\mu_p\)-torsor on the proper curve: locally, a zero or pole of
order \(pm\) is removed by rescaling \(W\) by the \(m\)-th power of a
local parameter. Its geometric monodromy is abelian of exponent \(p\),
so the torsor factors through the stipulated proper quotient. Pulling
it back along the restricted section over \(K\) gives \(\xi_i\).

For a finite place of \(K\) away from \(p\), the standard Fermat model
is smooth and proper, and the sections \(P_i\) are disjoint. The
divisor identity \eqref{eq:irreg-divisor} has no vertical component:
neither tangent line contains a fibre. The same local construction,
now with \(p\) invertible, extends the torsor to a finite etale torsor
over the integral model. Proper local points extend integrally, so
their evaluations are unramified. At a zero or pole this means
evaluation of the extended torsor, not of an undefined rational
function. A cusp section projects to the point section of its
proper centre, so it has the same property. Hence the valuation of
\(\xi_i\) is divisible by \(p\) away from \(\mathfrak p\).

At a first-case centre, \(a+b\equiv-c\not\equiv0\pmod p\).
Evaluation gives \eqref{eq:irreg-local-values}, and every displayed
value is a unit because \(\zeta_p\equiv1\pmod{\mathfrak p}\). Thus
\(\xi_i\in V\).

Finally, conjugating \(g_1\) by \(\sigma_i\) gives \(g_i\). The Kummer
isomorphism is natural when \(\mu_p\) carries its natural Galois
action. Pullback along the rational arithmetic section therefore
gives \(\xi_i=\sigma_i\xi_1\). Geometric conjugacy changes the
corresponding abelian cocycle by a coboundary only. There is no
identification with a constant coefficient module that could shift
the character.
\end{proof}

The same construction is available for each signed coordinate pair
on the original curve. No action of coordinate permutations on the
kernel of an arbitrary geometric quotient is required.

\subsection{A relative-class-rank criterion}
\label{sec:irreg-rank}

We give the local derivative argument in the precise form needed
here. Its classical antecedent is Eichler's method, as presented in
\cite[Section 3.2, pp.~107--108]{Gras2010}. The explicit sufficient
inequality below is used without an optimality assertion.

\begin{lemma}\label{lem:irreg-derivative}
Let \(F(T)\) be a rational function over \(\Z_p\), with a denominator
that is a unit at \(T=1\). Suppose \(F(1)=1\) and
\(F(\zeta_p)=b^p\) for \(b\in\mathcal O_L^\times\). Then
\begin{equation}
 \frac{F'(T)}{F(T)}=0
 \quad\text{in}\quad
 \F_p[T,\text{denominators}^{-1}]/(T-1)^p.
 \label{eq:irreg-derivative}
\end{equation}
\end{lemma}

\begin{proof}
Write \(b=B(\zeta_p)\), where \(B\in\Z_p[T]\) and
\(\deg B\leq p-2\). Clearing denominators and dividing by the monic
polynomial \(\Phi_p\) gives
\[
 F(T)-B(T)^p=\Phi_p(T)H(T)
\]
in a localization with denominator a unit at \(1\).
The residue of \(b\) is one, so \(B(1)\equiv1\pmod p\). Evaluation at
\(1\) gives
\[
 pH(1)=1-B(1)^p\equiv0\pmod{p^2}.
\]
Thus \(\overline H\) has a factor \(T-1\). Since
\(\Phi_p(T)\equiv(T-1)^{p-1}\pmod p\), we have
\(F-B^p=0\) modulo \((p,(T-1)^p)\). Differentiation preserves
\((T-1)^p\) in characteristic \(p\) and kills \(B^p\). The image of
\(F\) is a unit, giving \eqref{eq:irreg-derivative}.
\end{proof}

\begin{lemma}\label{lem:irreg-k-independent}
Suppose \(e,1-e,1-2e\in\Z_p^\times\), and let \(k\geq1\) satisfy
\begin{equation}
 k(k+1)-2<p.
 \label{eq:irreg-k-bound}
\end{equation}
Then the classes
\[
 d_i=
 \left[\frac{1-e+e\zeta_p^i}{1-e+e\zeta_p^{-i}}\right],
 \qquad 1\leq i\leq k,
\]
are linearly independent in \(L^\times/L^{\times p}\).
\end{lemma}

\begin{proof}
The bound implies \(k<p\). Suppose
\(\prod_i d_i^{\lambda_i}=1\), with
\(\lambda_i\in\{0,\ldots,p-1\}\) not all zero. The product of the
displayed unit representatives is therefore \(b^p\) for a unit
\(b\in\mathcal O_L^\times\). Apply
Lemma~\ref{lem:irreg-derivative} to
\[
 G(T)=\prod_{i=1}^k
 \left(\frac{1-e+eT^i}{1-e+eT^{-i}}\right)^{\lambda_i}.
\]
Its logarithmic derivative at \(1\) is
\(2e\sum_i i\lambda_i\), so \(\sum_i i\lambda_i\equiv0\pmod p\).
In particular \(\zeta_p^{\sum_i i\lambda_i}=1\) exactly. The identity
\[
 \frac{1-e+eT^i}{1-e+eT^{-i}}
   =T^i\frac{1-e+eT^i}{e+(1-e)T^i}
\]
then permits a second application to
\[
 F(T)=\prod_i
 \left(\frac{1-e+eT^i}{e+(1-e)T^i}\right)^{\lambda_i}.
\]
Writing
\(D_i(T)=(1-e+eT^i)(e+(1-e)T^i)\), we obtain
\[
 (2e-1)\sum_i\frac{\lambda_i iT^{i-1}}{D_i(T)}
       =0\pmod{p,(T-1)^p}.
\]
Let \(I\) be the nonempty support of the \(\lambda_i\).
All \(D_i(1)=1\), and \(2e-1\) is a unit, so
\[
 P(T)=
 \sum_{i\in I}\lambda_i iT^{i-1}
       \prod_{\substack{j\in I\\j\neq i}}D_j(T)
\]
is divisible by \((T-1)^p\) over \(\F_p\). On the other hand,
\[
 \deg P\leq2\sum_{i\in I}i-\min(I)-1
          \leq k(k+1)-2<p.
\]
Hence \(P=0\). For \(i_0=\min(I)\), its coefficient of \(T^{i_0-1}\)
is
\[
 \lambda_{i_0}i_0[e(1-e)]^{|I|-1}\neq0,
\]
and no other summand contributes at that degree. This contradiction
proves the lemma.
\end{proof}

\begin{theorem}\label{thm:irreg-rank}
Under the hypotheses of Section~\ref{sec:irreg-setup}, if
\begin{equation}
 (r^-+2)(r^-+3)-2<p,
 \qquad r^-=\dim_{\F_p}\Cl(K)[p]^-,
 \label{eq:irreg-rank-bound}
\end{equation}
then no permitted \(p\)-local centre is first case.
\end{theorem}

\begin{proof}
Assume otherwise and choose signed primitive unit coordinates
\(a,b,c\). Their nonzero residues satisfy \(a+b+c=0\), and, since
\(p\geq5\), are not all equal. Choose an unequal pair \(a,b\).
Then \(e=b/(a+b)\) satisfies the unit hypotheses of
Lemma~\ref{lem:irreg-k-independent}.

Lemma~\ref{lem:irreg-proper-torsors} gives
\(\xi_i\in V\), with \(\overline{\xi_i}=\xi_{-i}\). Thus
\(\xi_i/\xi_{-i}\in V^-\), and their localizations are the classes
\(d_i\). Take \(k=r^-+2\). By \eqref{eq:irreg-rank-bound} these \(k\)
local classes are independent; the global classes are therefore
independent as well. This contradicts
\(\dim_{\F_p}V^-=r^-+1\).
\end{proof}

\begin{corollary}\label{cor:irreg-residue}
Under \eqref{eq:irreg-rank-bound}, every permitted \(p\)-local
centre reduces to one of \(Q_0,Q_1,Q_\infty\).
\end{corollary}

\begin{proof}
For primitive integral coordinates the reduction of the equation is
\(X+Y=Z\). If they are not all units, exactly one is zero in the
residue field, giving one of the three displayed points. Moreover
\(\mu_p(\Q_p)=\{1\}\), so the only \(\Q_p\)-rational coordinate
boundary points themselves are the three rational cusps.
\end{proof}

\subsection{Odd-character Kummer congruences}
\label{sec:irreg-character}

In \(\F_p[E,W]/(W^{p-1})\), define
\begin{equation}
 C_j(E)=[W^j]\log(1-E+E\exp W),
 \qquad 1\leq j\leq p-2,
 \label{eq:irreg-C}
\end{equation}
where both formal series are truncated through degree \(p-2\).
All denominators are invertible in \(\F_p\).

\begin{theorem}\label{thm:irreg-characters}
Suppose that a section as in Section~\ref{sec:irreg-setup} has a
first-case \(p\)-local centre. For each ordered pair of its signed
coordinates, put \(e=b/(a+b)\bmod p\). For every odd integer
\(3\leq m\leq p-2\),
\begin{equation}
 \Cl(K)[p]^{(\omega^m)}=0
       \quad\Longrightarrow\quad C_m(e)=0.
 \label{eq:irreg-character-zero}
\end{equation}
In particular, without a class-component hypothesis,
\begin{equation}
 B_{p-m}\,C_m(e)=0\quad\text{in }\F_p.
 \label{eq:irreg-Bernoulli-congruence}
\end{equation}
\end{theorem}

\begin{proof}
By Lemmas~\ref{lem:irreg-virtual} and
\ref{lem:irreg-proper-torsors}, the vanishing class component implies
\(V^{(\omega^m)}=0\), and the section supplies
\(\xi_i=\sigma_i\xi_1\in V\). The direct \(\omega^m\)-projector is
\[
 e_m=\frac1{p-1}\sum_{i=1}^{p-1}i^{-m}\sigma_i.
\]
Choose integers \(c_i\) lifting its coefficients. The equality
\(e_m\xi_1=0\) means
\[
 \prod_{i=1}^{p-1}\xi_i^{c_i}=1
       \quad\text{in }K^\times/K^{\times p}.
\]
After localization at \(\mathfrak p\), it gives
\begin{equation}
 \prod_{i=1}^{p-1}(1-e+e\zeta_p^i)^{c_i}=\gamma^p
 \label{eq:irreg-projector-product}
\end{equation}
for a unit \(\gamma\in\mathcal O_L^\times\) with residue one.
Here and until reduction, \(e\) denotes its original value in
\(\Z_p\).

Set \(\lambda=\zeta_p-1\). Since \(v_\lambda(p)=p-1\) and
\(m+1\leq p-1\), there is an isomorphism
\begin{equation}
 \mathcal O_L/(\lambda^{m+1})
       \simeq R_m:=\F_p[t]/(t^{m+1}),\qquad
 \lambda\longmapsto t.
 \label{eq:irreg-trunc-ring}
\end{equation}
Indeed, \(\mathcal O_L=\Z_p[\lambda]\), and the Eisenstein equation
\(\Phi_p(1+\lambda)=0\) imposes no further relation among
\(1,\lambda,\ldots,\lambda^m\) in this quotient.
Writing \(\gamma=1+\lambda u\), the binomial formula shows that
\(\gamma^p\equiv1\pmod{\lambda^p}\). Thus the right side of
\eqref{eq:irreg-projector-product} equals \(1\) in \(R_m\).

Put
\[
 w=\log(1+t)=\sum_{r=1}^m\frac{(-1)^{r+1}t^r}{r}\in R_m.
\]
Truncated logarithm and exponential are inverse isomorphisms
between the multiplicative group \(1+(t)\) and the additive group
\((t)\). This follows by reducing their usual formal identities:
every denominator occurring through degree \(m\) has prime factors
less than \(p\). In particular \((1+t)^i=\exp(iw)\), and
\[
 \log\bigl(1-e+e(1+t)^i\bigr)
       =\sum_{j=1}^m C_j(e)i^j w^j.
\]
For \(1\leq j\leq m\),
\[
 \sum_{i=1}^{p-1}c_i i^j
       =\frac1{p-1}\sum_{i\in\F_p^\times}i^{j-m}
       =
       \begin{cases}
          0,&j<m,\\
          1,&j=m.
       \end{cases}
\]
The nonzero exponents \(j-m\) have absolute value at most \(p-3\),
so none is a multiple of \(p-1\). The logarithm of
\eqref{eq:irreg-projector-product} therefore reduces to
\(C_m(e)w^m=0\). Since \(w^m=t^m\neq0\), this proves
\eqref{eq:irreg-character-zero}.

For the Bernoulli assertion, use Herbrand's implication
\begin{equation}
 \Cl(K)[p]^{(\omega^m)}\neq0
       \quad\Longrightarrow\quad
       p\mid\operatorname{num}(B_{p-m}).
 \label{eq:irreg-Herbrand}
\end{equation}
The direct character convention in \eqref{eq:irreg-Herbrand} is that
of \cite[Theorem~1.1 and following paragraph, p.~151]{Ribet1976}; see also \cite[Section~6.3, pp.~100--102, Theorem~6.17]{Washington1982};
the projector there is written using \(\sigma_i^{-1}\) and agrees
with ours after reindexing. See also
\cite[Theorem 3.7 and the following paragraph, pp.~104--105]{Gras2010}.
If \(B_{p-m}\) is nonzero modulo \(p\), the contrapositive of
\eqref{eq:irreg-Herbrand} and
\eqref{eq:irreg-character-zero} give \(C_m(e)=0\).
Otherwise \eqref{eq:irreg-Bernoulli-congruence} is automatic.

The Bernoulli number has an unambiguous reduction: for
\(1\leq n\leq p-3\), its recurrence
\[
 \sum_{j=0}^n\binom{n+1}{j}B_j=0,\qquad B_0=1,
\]
shows inductively that \(B_n\in\Z_{(p)}\), since \(n+1<p\).
\end{proof}

\begin{remark}\label{rem:irreg-endpoint}
The proof includes \(m=p-2\). In that case
\(\mathcal O_L/(\lambda^{m+1})\) has characteristic \(p\) and length
\(p-1\), and no denominator divisible by \(p\) occurs. The logarithm
used above is a polynomial on a nilpotent ideal. It is not the
\(p\)-adic logarithm of \(\zeta_p\), which is zero. The reflected
character \(\omega^{1-m}\) in the usual notation for Mirimanoff
polynomials is likewise not the direct character of \(V\) used
here. The classical rational-point congruences are treated in
\cite[Theorems 3.5 and 3.7]{Gras2010}; the proper-torsor construction
above establishes their applicability under the stated quotient
hypotheses.
\end{remark}

\begin{corollary}\label{cor:irreg-cubic}
If \(\Cl(K)[p]^{(\omega^3)}=0\), no section under consideration has a
first-case \(p\)-local centre. In particular it suffices that
\begin{equation}
 p\nmid\operatorname{num}(B_{p-3}).
 \label{eq:irreg-cubic-condition}
\end{equation}
\end{corollary}

\begin{proof}
Expansion gives
\[
 C_1(E)=E,\qquad C_2(E)=\frac{E(1-E)}2,\qquad
 C_3(E)=\frac{E(1-E)(1-2E)}6.
\]
At a first-case centre, \(e\) and \(1-e\) are nonzero. Thus
Theorem~\ref{thm:irreg-characters} forces \(e=1/2\), or \(a=b\), for
each signed coordinate pair. All three nonzero signed residues
would be equal and have sum zero, impossible for \(p\geq5\).
Herbrand gives the stated sufficient Bernoulli condition.
\end{proof}

\begin{corollary}\label{cor:irreg-multiples}
Let \(n\) be an odd multiple of \(p\). The rank criterion, the
odd-character congruences, and the cubic corollary apply to sections
on \(U_n\) whose geometric quotient retains the pullbacks of the
proper mod-\(p\) abelian torsors on \(F_p\), with full \(G_\Q\) and
full local point-or-cusp representatives. In particular they apply
to its geometrically maximal prosolvable arithmetic fundamental
group.
\end{corollary}

\begin{proof}
The coordinate-power map
\[
 F_n\longrightarrow F_p,\qquad
 (X:Y:Z)\longmapsto
 (X^{n/p}:Y^{n/p}:Z^{n/p})
\]
takes boundary to boundary and first-case \(p\)-local points to
first-case points. Pullback retains the finite proper torsors used
in the proofs, and the full local representatives map to full
local representatives. Their geometric monodromy is abelian, so
these pullbacks factor through the maximal prosolvable quotient.
\end{proof}


The preceding rank criterion and character congruences, with the
cubic consequence, prove
Theorem~\ref{thm:first-case-obstructions-intro}.
